% !TeX root = ../../Book2.tex
\chapter{奇异积分与 Calderón–Zygmund 理论}
\label{b2:ch:18}
\BookChapterNumber{b2:ch:18}
\begin{chapterabstract}
若卷积核在原点附近像距离的负维数次幂那样增长，取绝对值后的积分通常发散，Young 不等式便不能直接使用。然而，核的正负部分可能相互抵消，使对称截断仍有极限。本章从 Hilbert 变换的具体计算出发，说明怎样把这种抵消转化为范数估计：先将输入分成有界的好部与零平均的坏块，再把坏块的零平均移到核的差分上。由此得到弱端点、中间指数的有界性与主值实现，最后用于共轭调和函数和椭圆方程的二阶导数估计。
\end{chapterabstract}

\begin{learninggoals}\begin{itemize}
\item\label{b2:item:18-goal-1}计算 Hilbert 变换的主值与 Fourier 符号，区分平方范数有界和点态极限。
\item\label{b2:item:18-goal-2}构造 Calderón--Zygmund 分解，说明坏块的零平均怎样改善远场衰减。
\item\label{b2:item:18-goal-3}核对标准核的大小、光滑性及离支集表示，并辨认平方有界性这一独立假设。
\item\label{b2:item:18-goal-4}完成弱型、插值、伴随与极大截断的论证，把稠密类上的主值传给一般输入。
\item\label{b2:item:18-goal-5}用 Riesz 变换描述 Poisson 共轭结构、频率投影与椭圆二阶估计。
\end{itemize}\end{learninggoals}

对区间示性函数作 Hilbert 变换，答案是一个对数。输入有界且可积，输出却在区间端点无界，并在无穷远只按倒数衰减。另一方面，它的 Fourier 符号模长为一，平方范数完全保持。同一个算子为什么能够同时具有这些性质？要回答这个问题，必须分清两个作用：绝对值估计控制大小，正负抵消则能消去绝对值估计中最危险的部分。

\ProseParagraphBreak

\chapref{b2:ch:17}已经准备了三件工具。二进微分定理让局部平均反映函数值，极大立方体分解把大平均的位置收集起来，Marcinkiewicz 插值则把弱端点与平方有界性接成中间的强型估计。本章调用它们时仍会写出所需版本及输入条件。Fourier 部分沿用\secref{b2:sec:0401}的 Plancherel 延拓，归一化固定为 \(e^{-2\pi\mathrm i x\cdot\xi}\)；涉及分布时，只使用\chapref{b2:ch:05}中的测试函数配对、分布导数和 Fourier 变换。除\subsecref{b2:subsec:180502}的周期比较外，底空间均为带 Lebesgue 测度的 \(\mathbb R^d\)。

\ProseParagraphBreak

首次阅读可先沿 Hilbert 变换、好坏分解、标准核、弱型证明这一顺序走到\subsecref{b2:subsec:180403}。这样先看清一个具体核怎样进入一般定理。随后读极大截断与主值传递，再到\secref{b2:sec:1805}看 Riesz 变换的实现和应用。向量值估计与 BMO 端点解释的是同一方法怎样服务于后续章节，可以在掌握标量证明后回读；它们不作为本章标量 \(L^p\) 有界性的前提。

\ProseParagraphBreak

进一步阅读可参考 Elias M. Stein 的《Singular Integrals and Differentiability Properties of Functions》\cite{Stein1970Singular}与 Loukas Grafakos 的《Classical Fourier Analysis》\cite{Grafakos2014Classical}。前者适合继续理解核抵消与微分性质的联系，后者可用于扩展奇异积分、乘子与函数空间的背景。本章采用的标准核版本在正文中证明，较弱核条件下的进一步结论另行说明出处和适用范围。

% !TeX root = ../../Book2.tex
\section{Hilbert 变换}
\label{b2:sec:1801}

在实直线上，核 \(1/x\) 的绝对积分在原点和无穷远都发散。然而它是奇函数，原点两侧的对称截断可以相互抵消。Hilbert 变换把这种抵消写成一个 Fourier 乘子，是研究奇异积分的第一个完整模型。

\subsection{主值}
\label{b2:subsec:180101}

对 \(f\in\mathcal S(\mathbb R)\)，定义\term[hilbert-bianhuan]{Hilbert 变换}[Hilbert transform]
\[
 Hf(x)\coloneqq\lim_{\varepsilon\downarrow0}
 \dfrac1\pi\int_{|x-y|>\varepsilon}\dfrac{f(y)}{x-y}\,\mathrm dy.
\]
\symidx{\(Hf\)：Hilbert 变换}
这里的极限称为\term[zhuzhi-jifen]{主值积分}[principal value integral]。换元 \(t=x-y\) 后，近原点的积分可以写成
\[
 \int_{\varepsilon<|t|<1}\dfrac{f(x-t)-f(x)}{t}\,\mathrm dt.
\]
均值定理使差商有界，所以极限存在；远处因 \(f\) 快速衰减而绝对可积。减去 \(f(x)\) 只在对称区间上进行，不能对一个无穷的常数积分作没有解释的减法。

\ProseParagraphBreak

若 \(f=\mathbf1_{(0,1)}\)，避开两个端点可直接计算
\[
 Hf(x)=\dfrac1\pi\log\left|\dfrac{x}{x-1}\right|.
\]
在区间内部，原点附近的对称抵消仍给出同一公式；在两个端点主值发散。这个例子同时显示主值不是绝对积分，以及有界输入可以产生无界输出。函数的单个点值改变不影响积分，但主值是否存在是对观察点而言的条件。

\subsection{Fourier 乘子}
\label{b2:subsec:180102}

考虑有限截断核 \(h_{\varepsilon,R}(t)=\mathbf1_{\{\varepsilon<|t|<R\}}/(\pi t)\)。其 Fourier 变换为
\[
 \widehat h_{\varepsilon,R}(\xi)
 =-\dfrac{2\mathrm i}{\pi}\int_\varepsilon^R\dfrac{\sin(2\pi\xi t)}t\,\mathrm dt.
\]
变换的符号由本卷负号指数决定。正、负频率得到相反的虚数值，不存在额外的 \(2\pi\) 幅值因子。
\begin{lemma}\label{b2:thm:18-dirichlet-integral}
上述乘子对 \(\varepsilon,R,\xi\) 一致有界，且在 \(\xi\ne0\) 时当 \(\varepsilon\downarrow0\)、\(R\uparrow\infty\) 收敛到 \(-\mathrm i\operatorname{sgn}\xi\)。
\end{lemma}
\begin{proof}
换元把积分化为 \(\int_a^b\sin u/u\,\mathrm du\)。在 \((0,1)\) 上被积函数的模不超过一；在 \((1,\infty)\) 上分部积分，使任意尾段的积分模不超过 \(3/a\)。这同时给出一致有界性和尾积分收敛。为计算全积分，先加入 \(e^{-\delta u}\)，对频率参数 \(s\) 求导得到
\[
 \dfrac{\mathrm d}{\mathrm ds}\int_0^\infty
 e^{-\delta u}\dfrac{\sin(su)}u\,\mathrm du
 =\int_0^\infty e^{-\delta u}\cos(su)\,\mathrm du
 =\dfrac{\delta}{\delta^2+s^2}.
\]
零参数处积分为零，故在 \(s=1\) 为 \(\arctan(1/\delta)\)。分部积分同样控制有阻尼和无阻尼的远尾，先固定尾端再令 \(\delta\downarrow0\)，得到 \(\pi/2\)。负频率由奇性得到负号。
\end{proof}
乘子在零点的取值不影响 \(L^2\) 算子；这与核在原点的截断约定是两个不同层次的无关零测修改。

\subsection{\texorpdfstring{\(L^2\)}{L2} 有界性}
\label{b2:subsec:180103}

\begin{theorem}\label{b2:thm:18-hilbert-l2}
Hilbert 变换具有唯一的 \(L^2(\mathbb R)\) 延拓，满足
\[
 \widehat{Hf}(\xi)=-\mathrm i\operatorname{sgn}(\xi)\widehat f(\xi),
 \qquad \|Hf\|_2=\|f\|_2.
\]
有限截断 \(h_{\varepsilon,R}*f\) 在 \(L^2\) 中趋于 \(Hf\)。
\end{theorem}
\begin{proof}
由 Plancherel 和上一引理，乘子的绝对值有共同上界，且几乎处处趋于绝对值为一的极限。对 \( |\widehat f|^2\) 使用支配收敛，得到截断在 \(L^2\) 中收敛及等距性。对 Schwartz 函数，主值积分已逐点存在，\(L^2\) 收敛抽出的几乎处处收敛子列识别两种定义。Schwartz 稠密性给出唯一延拓。
\end{proof}
因此 \(H^2f=-f\) 在 \(L^2\) 中成立，且 Hilbert 空间伴随为 \(-H\)。它不是正算子：不交区间之间的核可能为正也可能为负；其有界性来自 Fourier 符号和空间抵消，不能用正核的积分大小证明。

\subsection{\texorpdfstring{\(L^p\)}{Lp} 有界性}
\label{b2:subsec:180104}

核 \(K(x,y)=1/(\pi(x-y))\) 满足
\[
 |K(x,y)|\leq C|x-y|^{-1},\qquad
 |K(x,y)-K(x,y')|\leq C\dfrac{|y-y'|}{|x-y|^2}
\]
当 \(|y-y'|\leq|x-y|/2\) 时成立。第二项由差分分母直接计算。对于紧支集输入，在支集外不需要主值，\(Hf\) 等于普通核积分。这三项事实——\(L^2\) 有界、远离对角线的核表达、核差分控制——是稍后一般定理的全部输入。

\ProseParagraphBreak

由\cref{b2:thm:18-cz-lp}得到 \(1<p<\infty\) 的 \(L^p\) 有界性。这里是后文同章定理的应用位置，证明将在\secref{b2:sec:1804}完成。对任意 \(L^p\) 输入，连续延拓首先定义等价类；主值表示则需要控制截断的最大值，不能从强范数有界性独自推出。对本模型，这项控制在该节末尾给出。

\ProseParagraphBreak

无穷端点不能补上：区间示性函数的对数奇性已经否定 \(L^\infty\to L^\infty\)。在 \(L^1\) 端点，其远处渐近为 \(1/(\pi x)\)，也不绝对可积。后面的弱 \((1,1)\) 结论允许这样的尾部。

% 扩写定位：b2:subsec:180104
Hilbert 变换的两个端点都可以用一个函数看出障碍。对于 \(\mathbf1_{(0,1)}\)，已计算的主值是
\(\pi^{-1}\log|x/(x-1)|\)。在 \(x=0,1\) 附近，它具有对数无界性，故 Hilbert 变换不能把全部 \(L^\infty\) 输入送入 \(L^\infty\)。在 \(|x|\to\infty\) 时，展开 \(\log(1+1/(x-1))\) 得到 \(H\mathbf1_{(0,1)}(x)=1/(\pi x)+O(x^{-2})\)，所以其绝对值不可积，否定强型 \((1,1)\)。

\ProseParagraphBreak

两个失败来自不同的区域：无穷端点失败出现在输入跳跃的边缘，积分端点失败出现在远场。函数本身既可积又有界，不能通过要求输入同时属于两个端点来消去这些障碍。弱型 \((1,1)\) 仍然允许 \(1/x\) 的尾部，而 \(1<p<\infty\) 时对数奇性和 \(1/x\) 尾部都可以积到 \(p\) 次；这解释了强型区间的几何位置。

\ProseParagraphBreak

若输入同时具有零积分，远场的第一项会消失。例如取 \(f=\mathbf1_{(0,1)}-\mathbf1_{(1,2)}\)，则 \(Hf(x)=O(x^{-2})\) 于无穷处。一般紧支集 \(f\) 且 \(\int f=0\) 时，
\[
 Hf(x)=\frac1\pi\int f(y)
 \left(\frac1{x-y}-\frac1x\right)\mathrm dy,
\]
当 \(|x|\) 大于支集半径的两倍，绝对值不超过 \(C|x|^{-2}\int|y f(y)|\,\mathrm dy\)。这就是坏块尾估计的一维具体形式。不过零积分没有消除区间端点的对数奇性，仍不能恢复 \(L^\infty\) 强型。为讨论这类局部平均振荡，需要\appref{b2:app:C}的 BMO 空间。

% !TeX root = ../../Book2.tex
\section{Calderón–Zygmund 分解}
\label{b2:sec:1802}

核的奇性不能靠输入处处小来控制，因为一个 \(L^1\) 函数可能具有任意高的尖峰。Calderón–Zygmund 分解把尖峰限制在可计量的坏区域，并使每块坏函数的总质量为零。

\subsection{好/坏分解}
\label{b2:subsec:180201}

下面采用\term[calderon-zygmund-fenjie]{Calderón--Zygmund 分解}[Calderón--Zygmund decomposition]的下述版本。

\begin{theorem}[Calderón--Zygmund 分解]\label{b2:thm:18-cz-decomposition}
设 \(f\in L^1(\mathbb R^d)\)、\(\alpha>0\)。存在互不相交的二进立方体 \(Q_j\) 以及 \(f=g+\sum_jb_j\)，使
\[
 \sum_j|Q_j|\leq\alpha^{-1}\|f\|_1,\quad
 \|g\|_\infty\leq2^d\alpha,\quad \|g\|_1\leq\|f\|_1,
\]
\[
 b_j=(f-f_{Q_j})\mathbf1_{Q_j},\quad
 \int b_j=0,\quad
 \sum_j\|b_j\|_1\leq2\|f\|_1,
\]
其中 \(f_Q=|Q|^{-1}\int_Q f\)，级数在 \(L^1\) 中收敛。
\end{theorem}
\begin{proof}
对 \(\{M_{\mathcal D}f>\alpha\}\) 取\cref{b2:thm:17-dyadic-level}的极大立方体。定义 \(g=f\) 在其并的补集上，\(g=f_{Q_j}\) 在每个 \(Q_j\) 内。立方体平均上界和二进微分结论分别控制内、外的本性上确界。\(\int_{Q_j}|g|\leq\int_{Q_j}|f|\)，所以 \(L^1\) 界成立。坏块积分为零，且 \(\|b_j\|_1\leq2\int_{Q_j}|f|\)。不交性给出范数和与收敛，逐点除坏区域外全部项为零，坏区域内至多一项非零，故分解等式成立。
\end{proof}
复值 \(f\) 的平均仍是复数，控制来自 \(|f_Q|\leq(|f|)_Q\)，不需要把正部负部分别分解。如果原函数还属于 \(L^2\)，则 \(g\in L^2\)，而坏块级数也在 \(L^2\) 中收敛：其尾部平方积分等于不交支集上平方积分的尾和。

\ProseParagraphBreak

好函数的名称只表示它有受控的高度，并不意味着连续或光滑。坏函数虽保留尖峰，却拥有零平均；这个条件会在远场去掉核的常数项。

\subsection{弱型 \texorpdfstring{\((1,1)\)}{(1,1)}}
\label{b2:subsec:180202}

\begin{lemma}\label{b2:thm:18-cancellation-tail}
设 \(b\in L^1\) 支撑于边长 \(\ell\)、中心 \(c\) 的立方体 \(Q\)，且 \(\int b=0\)。若核满足
\[
 |K(x,y)-K(x,c)|\leq C\dfrac{|y-c|^\delta}{|x-c|^{d+\delta}}
 \quad(|x-c|\geq2|y-c|)
\]
其中 \(\delta>0\)，则在 \(Q^*=8\sqrt d\,Q\) 外有
\[
 \int_{\mathbb R^d\setminus Q^*}
 \left|\int_QK(x,y)b(y)\,\mathrm dy\right|\,\mathrm dx
 \leq C_{d,\delta}\|b\|_1.
\]
\end{lemma}
\begin{proof}
零平均使积分等于 \(\int_Q(K(x,y)-K(x,c))b(y)\,\mathrm dy\)。在扩大立方体外，\(|x-c|\geq4\sqrt d\,\ell\)，而 \(|y-c|\leq\sqrt d\,\ell/2\)，故差分假设适用。先对 \(y\) 取绝对积分，再对 \(x\) 用极坐标，尾核积分为
\[
 C\ell^\delta\int_{|x-c|\geq4\sqrt d\,\ell}|x-c|^{-d-\delta}\,\mathrm dx
 \leq C_{d,\delta}.
\]
Tonelli 完成积分交换。
\end{proof}
若不减去 \(K(x,c)\)，大小条件只有 \(|x-c|^{-d}\)，其远场绝对积分发散。这说明分解中的零平均不是方便的附属性质，而是端点估计中不可替代的衰减来源。扩大倍数的具体值只为了保证差分假设，后续常数允许依赖维数。

\subsection{插值}
\label{b2:subsec:180203}

将好部分交给 \(L^2\) 理论、坏部分交给上段尾估计，就能得到弱 \((1,1)\)。设暂时 \(f\in L^1\cap L^2\)，\(T\) 在 \(L^2\) 上有界，取 \(\alpha=t\)。好部分满足
\[
 |\{|Tg|>t/2\}|\leq4t^{-2}\|T\|_{2\to2}^2\|g\|_2^2
 \leq C t^{-1}\|f\|_1,
\]
因为 \(\|g\|_2^2\leq\|g\|_\infty\cdot\|g\|_1\)。坏立方体扩大后的并也有不超过 \(C t^{-1}\|f\|_1\) 的测度。在其补集上，核差分给出 \(\int|T\sum_jb_j|\leq C\sum_j\|b_j\|_1\)。对这项使用 \(L^1\) Chebyshev 即得同一界。

\ProseParagraphBreak

若坏块有无穷多个，先对有限和应用尾估计，再使用 \(L^2\) 收敛抽取几乎处处收敛子列和 Fatou。不能假定算子输出的无穷级数先在 \(L^1\) 中收敛；结论只在坏区域的补集上有该积分控制。下一节把核与算子的准确假设固定下来。

% !TeX root = ../../Book2.tex
\section{Calderón–Zygmund 核}
\label{b2:sec:1803}

Hilbert 核与 Riesz 核都有两种不同的性质：核本身在对角线附近不可积，移开一个变量后得到的核之差却衰减得更快。前一种性质决定积分为什么需要截断，后一种性质使坏块的零平均能够发挥作用。本节把这些条件分开写出，再说明怎样把核与一个已经有平方估计的算子联系起来。

\subsection{大小条件}
\label{b2:subsec:180301}

设整数 \(d\geq1\)，\(K:(\mathbb R^d\times\mathbb R^d)\setminus\{x=y\}\to\mathbb C\) 可测。若存在有限常数 \(C_0\geq0\)，使对每个 \(x\ne y\) 都有
\[
 |K(x,y)|\leq C_0|x-y|^{-d},
\]
则称 \(K\) 满足标准大小条件。半径为 \(r\) 的环带有 \(r^d\) 量级的体积，所以这个上界在每一倍增环带上只给出常数量级的绝对积分；对角线附近有无穷多个这样的环带，不能直接求和。

\ProseParagraphBreak

若 \(f\in L^2\) 支撑在紧集 \(F\) 内，则 \(f\in L^1\)。在与 \(F\) 有正距离的点 \(x\) 处，大小条件给出
\[
 \int_F|K(x,y)|\cdot|f(y)|\,\mathrm dy
 \leq C_0\cdot\operatorname{dist}(x,F)^{-d}\cdot\|f\|_1<\infty.
\]
因此支集外的核积分具有普通绝对积分的意义。这里同时用了输入的有限支集和观察点到支集的距离；支集内部的积分尚未由此定义。

\ProseParagraphBreak

卷积型核写成 \(K(x,y)=k(x-y)\)。例如取有界球面函数 \(\Omega\)，令 \(k(z)=\Omega(z/|z|)\cdot|z|^{-d}\)。若光滑输入在 \(x\) 附近恒等于一，则局部截断积分含有
\[
 \int_{\epsilon<|z|<1}k(z)\,\mathrm dz
 =\left(\int_{S^{d-1}}\Omega\,\mathrm d\sigma\right)\cdot
    \log\frac1\epsilon.
\]
球面平均不为零时，这一项发散。齐次次数 \(-d\) 给出大小条件，却没有提供抵消；球面上的奇函数平均为零，Hilbert 核就是一维的例子。

\subsection{光滑性条件}
\label{b2:subsec:180302}

坏块的零平均使 \(K(x,y)\) 被替换为 \(K(x,y)-K(x,c_Q)\)。因此需要控制的是移动输入变量 \(y\) 后的变化。对上述离开对角线的核，固定 \(0<\delta\leq1\)。标准光滑条件要求存在有限常数 \(C_1\geq0\)，使
\[
 \begin{aligned}
 |K(x,y)-K(x,y')|&\leq C_1\frac{|y-y'|^\delta}{|x-y|^{d+\delta}}
 &&\left(|y-y'|\leq\frac{|x-y|}2\right),\\
 |K(x,y)-K(x',y)|&\leq C_1\frac{|x-x'|^\delta}{|x-y|^{d+\delta}}
 &&\left(|x-x'|\leq\frac{|x-y|}2\right).
 \end{aligned}
\]
所有出现的原始点均满足 \(x\ne y\)。距离限制使新点仍离开对角线，且新旧距离可互相比较。第一行控制原算子的坏块；取伴随后输入、输出互换，第二行才成为伴随坏块所需的条件。

\ProseParagraphBreak

核的梯度界是一个常用的检查方法。若 \(K\) 在对角线外为 \(C^1\)，且
\[
 |\nabla_xK(x,y)|+|\nabla_yK(x,y)|\leq C|x-y|^{-d-1},
\]
则沿连接新旧点的线段积分导数即可取 \(\delta=1\)。例如 \(|y-y'|\leq|x-y|/2\) 时，线段上每点到 \(x\) 的距离至少为 \(|x-y|/2\)，所以积分中的梯度有同一个上界。

\ProseParagraphBreak

弱型证明实际上只用到了输入变量差的积分控制。设 \(h=|y-y'|>0\)，上述 Hölder 界给出
\[
 \begin{aligned}
 \int_{|x-y|>2h}|K(x,y)-K(x,y')|\,\mathrm dx
 &\leq C_1|S^{d-1}|h^\delta
            \int_{2h}^\infty r^{-1-\delta}\,\mathrm dr\\
 &=\frac{C_1|S^{d-1}|}{\delta\,2^\delta}.
 \end{aligned}
\]
\(h=0\) 时差为零。因而标准光滑性蕴含积分 Hörmander 条件
\[
 \sup_{y,y'}\int_{|x-y|>2|y-y'|}
       |K(x,y)-K(x,y')|\,\mathrm dx<\infty.
\]
它直接对应上一节对坏块远场的积分；点态条件还将用于截断比较和平均振荡估计。

\ProseParagraphBreak

也可以把幂次换成较慢的连续模量。设 \(\omega:[0,1]\to[0,\infty)\) 单调不减、\(\omega(0)=0\)，且
\[
 |K(x,y)-K(x,y')|\leq C|x-y|^{-d}\cdot
 \omega\!\left(\frac{|y-y'|}{|x-y|}\right)
 \quad\left(|y-y'|\leq\frac{|x-y|}2\right).
\]
径向积分换元 \(s=h/r\)，得到
\[
 \int_{|x-y|>2h}|K(x,y)-K(x,y')|\,\mathrm dx
 \leq C|S^{d-1}|\int_0^{1/2}\omega(s)\frac{\mathrm ds}s.
\]
因此 \(\int_0^1\omega(s)\,\mathrm ds/s<\infty\) 也足以控制坏块尾部。若已经给定有界 \(L^2\) 算子及支集外表示，这个积分条件可代替弱型证明中的输入变量 Hölder 条件；伴随也满足对应条件时，再由插值和对偶得到全部有限的中间指数。

\ProseParagraphBreak

取 \(\omega(s)=s^\delta\) 就回到幂次情形。取 \(\omega(s)=\bigl(\log(e/s)\bigr)^{-2}\) 时，换元 \(v=\log(e/s)\) 得到可积尾部 \(\int^\infty v^{-2}\,\mathrm dv\)；改成 \(\bigl(\log(e/s)\bigr)^{-1}\) 则发散。这里区分的是这一尾积分论证能否成立，并未排除某个具体算子还有别的抵消性质。按二进环带估计，同一个条件写成
\(\sum_{j\geq1}\omega(2^{-j})<\infty\)：每个环带的体积消去核的 \(|x-y|^{-d}\)，留下的正是模量。单调性给出
\[
 (\log2)\cdot\omega(2^{-j-1})
 \leq\int_{2^{-j-1}}^{2^{-j}}\omega(s)\frac{\mathrm ds}s
 \leq(\log2)\cdot\omega(2^{-j}),
\]
从而级数与积分具有相同的收敛性。

\subsection{奇异积分算子}
\label{b2:subsec:180303}

核的条件只描述离开对角线的作用。还须指定一个算子，并独立证明其平方有界性。
\Needspace{6\baselineskip}
\begin{definition}\label{b2:thm:18-cz-operator-definition}
设整数 \(d\geq1\)，\(T:L^2(\mathbb R^d;\mathbb C)\to L^2(\mathbb R^d;\mathbb C)\) 为有界复线性算子。设可测核 \(K:(\mathbb R^d\times\mathbb R^d)\setminus\{x=y\}\to\mathbb C\)，并存在有限常数 \(C_0,C_1\geq0\)、\(0<\delta\leq1\)，使
\[
 |K(x,y)|\leq C_0|x-y|^{-d}
 \quad(x\ne y),
\]
\[
 \begin{aligned}
 &|K(x,y)-K(x',y)|+|K(y,x)-K(y,x')|\\
 &\quad\leq C_1\frac{|x-x'|^\delta}{|x-y|^{d+\delta}}
 \quad\left(x\ne y,\ |x-x'|\leq\frac{|x-y|}2\right).
 \end{aligned}
\]
若对每个支撑在某个紧集 \(F\) 内的 \(f\in L^2(\mathbb R^d)\)，在 \(\mathbb R^d\setminus F\) 上几乎处处有
\[
 Tf(x)=\int_FK(x,y)\cdot f(y)\,\mathrm dy,
\]
则称 \(T\) 为具有核 \(K\) 的\term[calderon-zygmund-suanzi]{Calderón--Zygmund 算子}[Calderón--Zygmund operator]。上述核条件称为标准核条件；第二个显示式同时包含两个变量的光滑性。
\end{definition}

平方有界性没有被核条件包含。例如零核既对应零算子，也对应任意有界乘法算子 \(M_af=af\)、\(a\in L^\infty\)：输入支集外两种输出都为零。后面的有界性理论适用于这些算子，但若只知道支集外的核，就不能判定给定算子是否等于硬截断的主值。

\ProseParagraphBreak

伴随所需的核可以直接从内积求出。取紧支集 \(f,g\in L^2\)，并假设两个支集有正距离。它们均可积，且
\[
 \int\!\int |K(x,y)|\cdot|f(y)|\cdot|g(x)|\,\mathrm dy\,\mathrm dx
 \leq C\|f\|_1\cdot\|g\|_1<\infty.
\]
这里 \(C\) 还依赖两个支集的距离。Fubini 与 \(\langle Tf,g\rangle=\langle f,T^*g\rangle\) 给出
\[
 \int f(y)\cdot\overline{T^*g(y)}\,\mathrm dy
 =\int f(y)\cdot
       \overline{\int\overline{K(x,y)}\cdot g(x)\,\mathrm dx}\,\mathrm dy.
\]
固定 \(g\)，让 \(f\) 遍历任一离开 \(\operatorname{supp}g\) 的有界紧区域上的 \(L^2\) 函数，即可识别两边的局部 \(L^2\) 函数。再以可数个这样的区域穷竭支集外，得到伴随的核
\[
 K^*(x,y)=\overline{K(y,x)}.
\]
它的大小与两个变量光滑常数仍可取 \(C_0,C_1\)，且 \(\|T^*\|_{2\to2}=\|T\|_{2\to2}\)。输入变量的光滑性在伴随后变成输出变量的光滑性，两者不能只保留其中一条。

\ProseParagraphBreak

下面从一个具体核构造算子。这里的目标是先取得独立的 \(L^2\) 界，再使用下一节的一般有界性定理。
\Needspace{6\baselineskip}
\begin{theorem}\label{b2:thm:18-homogeneous-kernel}
设整数 \(d\geq2\)，\(\Omega\in C^1(S^{d-1};\mathbb C)\)，且 \(\int_{S^{d-1}}\Omega\,\mathrm d\sigma=0\)，其中 \(\sigma\) 为球面面积测度。则对每个 \(f\in\mathcal S(\mathbb R^d)\)、每个 \(x\in\mathbb R^d\)，主值
\[
 T_\Omega f(x)=\lim_{\epsilon\downarrow0}
 \int_{|y|>\epsilon}\frac{\Omega(y/|y|)}{|y|^d}\cdot f(x-y)\,\mathrm dy
\]
存在。它给出兼容的有界线性算子 \(T_\Omega:L^p(\mathbb R^d)\to L^p(\mathbb R^d)\)、\(1<p<\infty\)，且
\[
 \|T_\Omega f\|_p\leq C_{d,p}\|\Omega\|_{C^1(S^{d-1})}\cdot\|f\|_p.
\]
其 \(L^2\) Fourier 乘子在 \(\xi\ne0\) 时为
\[
 m(\xi)=-\int_{S^{d-1}}\Omega(\theta)\cdot
 \left(\log|2\pi\xi\cdot\theta|
       +\frac{\pi\mathrm i}{2}\cdot\operatorname{sgn}(\xi\cdot\theta)\right)\mathrm d\sigma(\theta).
\]
对数奇性按球面上的可积函数解释，\(m\) 为零次齐次；其在 \(\xi=0\) 的取值可任意指定。
\end{theorem}
\begin{proof}
先固定 \(f\in\mathcal S\)、\(x\)。在 \(0<|y|<1\) 的积分中减去 \(f(x)\)，球面平均零使常数部分消失。由微积分基本定理，
\[
 |f(x-y)-f(x)|\leq |y|\sup_{|z-x|\leq1}|\nabla f(z)|,
\]
所以近原点的绝对积分由 \(C_x\int_0^1\mathrm dr\) 控制。远处 \(|x-y|\) 随 \(|y|\) 增大，Schwartz 衰减使径向尾部可积。因此主值在每个点存在。

为了求平方界，先在径向乘上 \(e^{-\eta r}\)、\(\eta>0\)，得到同样按球面平均零定义的算子 \(T_\eta\)。对 \(\operatorname{Re}z>0\)，有
\[
 \int_0^\infty\frac{e^{-zr}-e^{-r}}r\,\mathrm dr=-\log z.
\]
右半平面取连续对数：在其每个紧子集上，对 \(z\) 求导的积分由一个指数衰减函数控制，导数为 \(-1/z\)；再由 \(z=1\) 的值为零确定常数。减项 \(e^{-r}\) 与角变量无关，乘 \(\Omega\) 后的积分为零。由此得到
\[
 m_\eta(\xi)=-\int_{S^{d-1}}\Omega(\theta)\cdot
                  \log(\eta+2\pi\mathrm i\xi\cdot\theta)\,\mathrm d\sigma(\theta).
\]
具体地，先在两个有限径向截断之间计算，再与 Schwartz 频率测试函数配对。近 \(r=0\) 时，差
\(e^{-(\eta+2\pi\mathrm i\xi\cdot\theta)r}-e^{-r}\)
除以 \(r\) 后由 \(C_\eta(1+|\xi|)\) 控制；乘测试函数后可积。远处则由 \(e^{-\eta r}+e^{-r}\) 控制。因而可以在这些配对中换序并放开截断，得到 \(\widehat{T_\eta f}=m_\eta\cdot\widehat f\) 的温和分布恒等式。

需要的估计是对 \(\eta\) 与频率同时一致的。对 \(\xi\ne0\)，记 \(\omega=\xi/|\xi|\)、\(a=\eta/|\xi|\)。平均零先去掉 \(\log|\xi|\)。若 \(a\geq1\)，再去掉 \(\log a\)，余下的 \(\log(1+2\pi\mathrm i\omega\cdot\theta/a)\) 一致有界。若 \(0<a\leq1\)，对数虚部的绝对值不超过 \(\pi/2\)，实部的正部一致有界，负部则由
\(C+\bigl|\log|\omega\cdot\theta|\bigr|\) 控制。旋转不变性给出
\[
 \sup_{\omega\in S^{d-1}}
 \int_{S^{d-1}}\bigl|\log|\omega\cdot\theta|\bigr|\,\mathrm d\sigma(\theta)<\infty:
\]
旋转到一个坐标方向后，在赤道附近令 \(s=\omega\cdot\theta\)，面积密度有界，该部分归结为
\(\int_{-1/2}^{1/2}\bigl|\log|s|\bigr|\,\mathrm ds<\infty\)；其余球面上对数有界。因此
\[
 \sup_{\eta>0}\|m_\eta\|_\infty\leq C_d\|\Omega\|_\infty.
\]
对几乎每个球面点令 \(\eta\downarrow0\)，复对数趋向
\(\log|2\pi\xi\cdot\theta|+\mathrm i\pi\cdot\operatorname{sgn}(\xi\cdot\theta)/2\)。上述控制给出 \(m_\eta(\xi)\to m(\xi)\)。Plancherel 与频率上的支配收敛于是给出有界 \(L^2\) 乘子算子 \(T\)，且 \(T_\eta f\to Tf\) 于 \(L^2\)。

在物理空间中，减去 \(f(x)\) 后的近原点积分和 Schwartz 远尾，都有与 \(0<\eta\leq1\) 无关的可积控制，所以 \(T_\eta f(x)\to T_\Omega f(x)\) 对每个 \(x\) 成立。取 \(L^2\) 收敛的几乎处处子列，便把这个主值识别为 \(Tf\)。平均零也从乘子公式中消去 \(\log|\xi|\)，给出零次齐次性。

核的大小与梯度常数由 \(C_d\|\Omega\|_{C^1}\) 控制。它对光滑紧支集输入的支集外表示已由主值公式给出；对一般紧支集 \(L^2\) 输入，在任意到其支集有正距离的有界区域上，用共用一个稍大紧支集的光滑 \(L^2\) 逼近。该区域到共同支集有正距离，核积分的差由 \(L^1\) 误差一致控制；\(L^2\) 输出收敛再抽子列，即得所需支集外表示。以可数个这样的区域穷竭支集外，说明 \(T\) 满足标准算子定义。\cref{b2:thm:18-cz-lp}因此给出全部 \(1<p<\infty\) 的界。先将非零的 \(\Omega\) 除以其 \(C^1\) 范数应用该结论，再还原尺度，就得到陈述中的线性范数因子；\(\Omega=0\) 时直接成立。
\end{proof}

乘子公式的两部分来自不同的对称性。若 \(\Omega\) 为奇函数，对数项的球面积分为零；若 \(\Omega\) 为偶函数，符号函数项的积分为零。再假设 \(\Omega\) 实值时，前一种乘子为纯虚函数，后一种为实偶函数。二维的 \(\Omega(\theta)=\theta_1\theta_2\) 给出偶核 \(x_1x_2/|x|^4\)，而 \(\Omega(\theta)=\theta_j\) 给出 Riesz 型奇核。相反，\(\Omega=1\) 产生前面的径向对数发散；人为扣除发散常数属于另外一种归一化，不能据此声称原硬截断主值已经存在。

\subsection{截断}
\label{b2:subsec:180304}

取整数 \(d\geq1\)，设核在对角线外为 Borel 可测，且满足 \(|K(x,y)|\leq C_0|x-y|^{-d}\)。前面的标准核满足这一可测性：在远离对角线的小邻域内，先移动 \(x\) 再移动 \(y\)，两个 Hölder 界共同给出连续性。对 \(\epsilon>0\) 定义硬截断
\[
 T_\epsilon f(x)=\int_{|x-y|>\epsilon}K(x,y)\cdot f(y)\,\mathrm dy.
\]
若 \(f\in L^p(\mathbb R^d)\)、\(1<p<\infty\)，Hölder 给出
\[
 \int_{|x-y|>\epsilon}|K(x,y)|\cdot|f(y)|\,\mathrm dy
 \leq C_{d,p}C_0\epsilon^{-d/p}\|f\|_p.
\]
这里径向控制核属于 \(L^{p'}\)，因为 \(dp'>d\)。\(p=1\) 时改用 \(|K|\leq C_0\epsilon^{-d}\)，仍有绝对收敛。故每个固定截断在这些输入上有普通积分的意义；一般有界输入却可能在无穷远处失去可积性。

\ProseParagraphBreak

把截断半径移向零，还需另外的论证。给定的标准 \(L^2\) 算子可以含有一个支集外核看不见的乘法项，所以截断的极限不必等于它。常见的识别办法，是先对 \(C_c^\infty\) 输入验证截断的逐点极限，再证明这些截断在 \(L^2\) 中趋于给定算子；平方收敛的子列才能把两个极限识别起来。下一节先证明极大截断的有界性，再把这项识别从光滑类传给一般 \(L^p\) 输入。

\ProseParagraphBreak

不同截断约定的比较也应写出实际差值。取实函数 \(\chi\in C^\infty([0,\infty))\)，在 \([0,1]\) 为零、在 \([2,\infty)\) 为一，且 \(0\leq\chi\leq1\)，令
\[
 T^\chi_\epsilon f(x)=\int K(x,y)\cdot
                 \chi(|x-y|/\epsilon)\cdot f(y)\,\mathrm dy.
\]
与硬截断的差仅在 \(\epsilon<|x-y|<2\epsilon\)，所以对有限指数的 \(L^p\) 输入，
\[
 |T^\chi_\epsilon f(x)-T_\epsilon f(x)|
 \leq C_dC_0\epsilon^{-d}\int_{B(x,2\epsilon)}|f|
 \leq C_d'C_0\cdot Mf(x).
\]
此界能传递 \(1<p<\infty\) 的极大范数估计，但右端不随 \(\epsilon\) 消失，不能单凭这一界识别极限。

\ProseParagraphBreak

若硬截断在某点已有有限极限，则两种截断其实有相同极限，而不必再要求核的抵消。由 \(\int_1^2\chi'(s)\,\mathrm ds=1\)，对上述绝对收敛积分换序，得
\[
 T^\chi_\epsilon f(x)=\int_1^2\chi'(s)\cdot T_{s\epsilon}f(x)\,\mathrm ds.
\]
因此它与硬截断极限 \(L\) 的差不超过
\[
 \left(\int_1^2|\chi'(s)|\,\mathrm ds\right)
 \sup_{\epsilon\leq r\leq2\epsilon}|T_rf(x)-L|\longrightarrow0.
\]
对球面平均零的齐次核，还能在光滑输入上直接比较：环带差核的积分为零，减去 \(f(x)\) 后，差不超过
\(C_d\|\Omega\|_\infty\epsilon\sup_{B(x,2\epsilon)}|\nabla f|\)。这给出共同的光滑类极限；向一般输入推广仍须控制极大算子的逼近误差。这里比较的是两种截断；它们与预先给定的算子 \(T\) 是否相等，仍是前一段所说的识别问题。

\ProseParagraphBreak

截断还可以先保留一个有限的外半径，再逐次放开。对 \(L^p\)、\(1\leq p<\infty\)，固定内半径后的远场已有上述绝对积分控制，外半径趋于无穷不会改变结果。对常函数或一般有界输入，远场可能发散或仅有条件收敛，需要另行指定归一化；有限指数上建立的定义不能不加说明地用于这些输入。

% !TeX root = ../../Book2.tex
\section{\texorpdfstring{\(L^p\)}{Lᵖ} 有界性}
\label{b2:sec:1804}

分解给出弱型，插值给出下半段指数，伴随再给出上半段。三个步骤各使用不同的假设；把它们连起来才得到完整的 \(L^p\) 范围。

\subsection{\texorpdfstring{\(L^2\)}{L2} 假设}
\label{b2:subsec:180401}

全章的主要假设是 \(L^2\) 有界性。对卷积算子，它通常由有界 Fourier 乘子得到；对非平移不变的算子，需要另外的估计或测试条件。标准核条件只控制远场差分，不能单独排除局部能量增长。

\ProseParagraphBreak

在分解中，\(L^2\) 只用于好函数：其高度有限，且 \(L^1\) 范数有限，所以它落在 \(L^2\)。这样将原输入的 \(L^1\) 大小与平方估计连接起来。坏函数并没有更好的平方范数，若把整个 \(f\) 直接代入 \(L^2\) 估计，会失去所需的端点常数。

\subsection{弱型 \texorpdfstring{\((1,1)\)}{(1,1)}}
\label{b2:subsec:180402}

\begin{theorem}[Calderón--Zygmund 弱型]\label{b2:thm:18-cz-weak}
设 \(T\) 满足\cref{b2:thm:18-cz-operator-definition}。则对 \(f\in L^1\cap L^2\)、\(t>0\)，有
\[
 |\{|Tf|>t\}|\leq C t^{-1}\|f\|_1,
\]
其中常数只依赖维数、光滑指数、核常数和 \(L^2\) 算子范数。\(T\) 有唯一的 \(L^1\to L^{1,\infty}\) 连续延拓。
\end{theorem}
\begin{proof}
取高度 \(t\) 的分解。好部分、扩大坏区域及补集上的坏部分估计已经在\secref{b2:sec:1802}完成，其和给出结论。对一般 \(L^1\) 输入取 \(f_n\in L^1\cap L^2\) 范数逼近。弱估计用于差 \(f_n-f_m\)，使 \(Tf_n\) 在每个有限测度集合上依测度 Cauchy。通过选择差的超水平集合测度可求和的子列，在整数半径球上同时得到几乎处处极限；这些局部极限在交集上一致，拼成可测函数 \(Tf\)。Fatou 的集合版本把弱界传给极限。对任意另一逼近列，同样的差估计使极限相等几乎处处，因此延拓唯一。这里没有把弱拟范数误当成满足常数一的三角不等式。
\end{proof}
弱端点一般不能升级为 \(L^1\) 有界，Hilbert 变换已给反例。下一步只在开区间内插值。

\subsection{\texorpdfstring{\(1<p<\infty\)}{1<p<∞}}
\label{b2:subsec:180403}

\begin{theorem}[Calderón--Zygmund]\label{b2:thm:18-cz-lp}
设 \(T\) 满足\cref{b2:thm:18-cz-operator-definition}。则对所有 \(1<p<\infty\)，有唯一的有界延拓 \(T:L^p\to L^p\)，且
\[
 \|Tf\|_p\leq C_p\|f\|_p.
\]
\end{theorem}
\begin{proof}
在 \(1<p<2\) 时，弱 \((1,1)\) 与强 \((2,2)\) 经\cref{b2:thm:17-marcinkiewicz}得到所需界。\(T^*\) 也是标准核算子，所以同样在 \(1<p'<2\) 有界。对 \(p>2\) 的 \(f\in L^p\cap L^2\) 与测试 \(g\in L^{p'}\cap L^2\)，
\[
 \left|\int Tf\,\overline g\right|
 =\left|\int f\,\overline{T^*g}\right|
 \leq C_{p'}\|f\|_p\cdot\|g\|_{p'}.
\]
对测试 \(g\) 取上确界得到 \(L^p\) 界；稠密性给出唯一延拓。公共输入上的各个延拓由 \(L^2\) 或局部依测度极限识别，因此它们是一个兼容算子的不同空间实现。
\end{proof}
核是否齐次、算子是否卷积都没有进入这项证明。其可迁移的部分是大小、光滑、局部平方理论和零平均分解，而不是某一个特定 Fourier 符号。

\subsection{对偶性}
\label{b2:subsec:180404}

最后说明主值的点态问题。若 \(T_\varepsilon\) 在 \(L^2\) 中实现 \(T\)，定义 \(T_*f=\sup_{\varepsilon>0}|T_\varepsilon f|\)。对 \(1<q<\infty\)，核差分和 \(T\) 的 \(L^q\) 有界性给出
\begin{proposition}\label{b2:thm:18-maximal-truncation}
在上述主值实现条件下，
\[
 T_*f(x)\leq C_q\left[
 M(|Tf|^q)(x)^{1/q}+M(|f|^q)(x)^{1/q}\right]
\]
几乎处处成立。因此 \(T_*\) 在每个 \(L^p\)、\(1<p<\infty\) 上有界。
\end{proposition}
\begin{proof}
固定 \(\varepsilon\)，写 \(f=f_1+f_2\)，其中 \(f_1=f\mathbf1_{B(x,2\varepsilon)}\)。对 \(z\in B(x,\varepsilon/4)\)，\(Tf_2(z)\) 是支集外的核积分。比较它与 \(T_\varepsilon f(x)\)，中间环带的积分由大小条件控制为 \(C Mf(x)\)；远场核差分分成半径 \(2^k\varepsilon\) 的环带，每带不超过 \(C2^{-k\delta}Mf(x)\)，可求和。故
\[
 |T_\varepsilon f(x)|\leq |Tf(z)|+|Tf_1(z)|+CMf(x).
\]
对 \(z\) 取 \(q\) 次平均并用 Minkowski。第一项被 \(M(|Tf|^q)^{1/q}\) 控制；第二项由 \(L^q\) 有界性不超过
\[
 C\varepsilon^{-d/q}\|f\mathbf1_{B(x,2\varepsilon)}\|_q
 \leq C M(|f|^q)(x)^{1/q}.
\]
用 Hölder 将 \(Mf\) 吸收到同一项，再对截断半径取上确界。对给定 \(p>1\)，选 \(1<q<p\)，极大算子在 \(L^{p/q}\) 上有界，而 \(T\) 在 \(L^p\) 上有界，得到结论。
\end{proof}
若某个稠密测试类上截断逐点收敛，以上最大估计使截断差的振荡集合由逼近误差控制，按\cref{b2:thm:17-differentiation}的相同误差集合方法得到一般 \(L^p\) 输入几乎处处收敛；范数极限识别它为 \(Tf\)。这适用于 Hilbert 变换。弱 \(L^1\) 的最大截断理论需要额外端点论证，留在\appref{b2:app:C}讨论。

% 扩写定位：b2:subsec:180404
后面的平方函数需要一次处理许多输出。逐个标量估计再求平方和，通常会产生与输出个数有关的常数，因而不足以处理无限频带。
\begin{proposition}\label{b2:prop:18-hilbert-valued-cz}
设 \(\mathcal H_1,\mathcal H_2\) 为有限维 Hilbert 空间。令 \(T:L^2(\mathbb R^d;\mathcal H_1)\to L^2(\mathbb R^d;\mathcal H_2)\) 有界，核 \(K(x,y)\) 为两空间之间的线性映射，满足统一的算子范数大小与两变量 Hölder 条件，并在支集外表示 \(T\)。则 \(T\) 满足弱型 \((1,1)\) 和所有 \(1<p<\infty\) 的强型估计；常数只依赖上述界、\(p,d,\delta\)，不依赖 Hilbert 空间维数。
\end{proposition}
\begin{proof}
对标量函数 \(\|f(x)\|_{\mathcal H_1}\) 作二进坏立方体选择。好函数在每个坏块内取 Bochner 平均 \(f_Q\)，坏函数为 \(b_Q=(f-f_Q)\mathbf1_Q\)。凸性给出
\(\|g\|_\infty\leq2^d\lambda\)、\(\|g\|_1\leq\|f\|_1\)，从而
\(\|g\|_2^2\leq2^d\lambda\|f\|_1\)。
且 \(\int b_Q=0\)、\(\sum_Q\|b_Q\|_1\leq2\|f\|_1\)，坏块的测度总和不超过 \(\|f\|_1/\lambda\)。

在放大的坏块外，把积分写成
\[
 Tb_Q(x)=\int_Q[K(x,y)-K(x,c_Q)]b_Q(y)\,\mathrm dy.
\]
先取 Hilbert 范数再用算子范数控制核差，得到
\(\int_{\mathbb R^d\setminus Q^*}\|Tb_Q(x)\|\,\mathrm dx\leq C\|b_Q\|_1\)。
好部分用 \(L^2\) Chebyshev，坏部分用刚才的积分界和 \(L^1\) Chebyshev，便重复得到弱型估计。每一步常数只来自标量几何与给定的算子范数界。

对 \(1<p<2\)，将输出的 Hilbert 范数作为次线性输出；输入按 \(\|f(x)\|\) 的高低截断，本章之前的插值证明仍逐行成立。对 \(p>2\)，在 Hilbert 对偶配对中使用伴随算子；其核为 \(K(y,x)^*\)，具有同一组界。因而得到全部指数。
\end{proof}
例如，若 \(T_j\) 是一族具有统一常数的卷积型 Calderón--Zygmund 算子，则对任意有限族 \(f_j\)，
\[
 \left\|\left(\sum_j|T_jf_j|^2\right)^{1/2}\right\|_p
 \leq C_p\left\|\left(\sum_j|f_j|^2\right)^{1/2}\right\|_p.
\]
取对角算子 \(T(f_j)_j=(T_jf_j)_j\)：其核的算子范数为 \(\sup_j|k_j(x-y)|\)，其 \(L^2\) 界由逐项 Plancherel 或逐项 \(L^2\) 界相加得到。有限维一致常数和单调收敛允许推广到可数族。这里只要求对角核统一；一般矩阵核不能以每个矩阵元有界代替算子范数有界，因为矩阵大小会影响其范数。

\ProseParagraphBreak

这一证明也适用于坏块估计中只给出统一积分 Hörmander 条件的情形，前提是算子及伴随都满足对应条件。最大截断的 Cotlar 估计仍需另外检查点态条件或其他替代工具；抽象 \(L^p\) 有界性与主值几乎处收敛是两件不同的事。

无穷输入端点也存在一个可完整验证的替代估计。对局部可积 \(h\)，记
\(\|h\|_{\mathrm{BMO}}=\sup_Q|Q|^{-1}\int_Q|h-h_Q|\)，其中上确界取所有立方体，且 \(h_Q\) 为平均。该空间将在\appref{b2:app:C}系统定义；这里先证明与本章核条件直接相关的结论。
\begin{proposition}\label{b2:prop:18-infinity-bmo}
标准 Calderón--Zygmund 算子 \(T\) 对每个有界紧支集 \(f\) 满足
\(\|Tf\|_{\mathrm{BMO}}\leq C\|f\|_\infty\)，常数与输入支集大小无关。
\end{proposition}
\begin{proof}
固定立方体 \(Q\)、中心 \(c_Q\)，取充分大的固定扩大块 \(Q^*\)。分解 \(f=f_1+f_2\)，其中 \(f_1=f\mathbf1_{Q^*}\)。局部部分由平方界控制：
\[
 \frac1{|Q|}\int_Q|Tf_1|
 \leq |Q|^{-1/2}\|T\|_{2\to2}\|f_1\|_2
 \leq C\|f\|_\infty.
\]
远场在 \(Q\) 内具有支集外核表示。取比较常数
\(c=\int_{(Q^*)^c}K(c_Q,y)f(y)\,\mathrm dy\)，该积分因紧支集及离开对角线而有限。核的第一变量光滑性给出
\[
 |Tf_2(x)-c|
 \leq C\|f\|_\infty\ell(Q)^\delta
 \int_{|y-c_Q|>c_d\ell(Q)}|y-c_Q|^{-d-\delta}\,\mathrm dy
 \leq C'\|f\|_\infty.
\]
故 \(Q\) 上相对于常数 \(c\) 的平均振荡受控。换成真正平均 \((Tf)_Q\) 至多增加因子二，取全部 \(Q\) 的上确界即可。
\end{proof}
两个核变量的职责在此再次分开：弱积分端点的坏块需要输入变量的差分，无穷端点的局部振荡需要输出变量的差分。仅给其中一种光滑性，不能把另一方向的结论也视为已经证明。

\ProseParagraphBreak

当输入仅有界而不紧支集时，上面常数 \(c\) 的远场积分可能发散。不过差
\(\int[K(x,y)-K(c_Q,y)]f(y)\,\mathrm dy\) 仍绝对收敛，因其多一个衰减幂。由此可以按模常数构造局部相容的输出；附录中的 BMO 对偶也自然只见到模常数的对象。这个归一化说明了为什么端点的正确目标是平均振荡，而不是一个具有绝对数值的全局有界函数。

% !TeX root = ../../Book2.tex
\OptionalSection{Schauder 估计}{b2:sec:1605}

待撰写。

\subsection{Hölder 空间与缩放}
\label{b2:subsec:160501}

待撰写。

\subsection{Poisson 方程的内部估计}
\label{b2:subsec:160502}

待撰写。

\subsection{冻结系数与一般内部估计}
\label{b2:subsec:160503}

待撰写。

\subsection{边界估计的条件与理论边界}
\label{b2:subsec:160504}

待撰写。


\begin{chaptersummary}
Calderón--Zygmund 分解在每个坏立方体上保留平均值，把剩余部分变成零平均函数。远离这个立方体时，核可以减去中心值；核差分比核本身多衰减一个正幂次，于是坏部的远场可以绝对积分。好部由平方估计控制，坏区由体积估计控制，这三项合起来得到弱 \((1,1)\)。再经插值与伴随，标准 Calderón--Zygmund 算子在全部 \(1<p<\infty\) 上有兼容的延拓。

\ProseParagraphBreak

使用这一结论时，还要分别核对核条件、\(L^2\) 有界性和截断实现。离支集核并不决定算子在对角线上的作用，\(L^p\) 范数延拓也不自动等于逐点主值。极大截断提供统一的误差控制，具体核在光滑稠密类上的收敛提供起点，两者结合才把主值传给一般输入。Hilbert 与 Riesz 变换满足这些条件，它们的零阶 Fourier 符号又把共轭调和结构和二阶椭圆估计连接起来。区域问题仍须经过截断和局部化，所得估计中的低阶项不能在没有额外条件时删去。
\end{chaptersummary}

\BookExercises


% 自拟：由区间模型精算分布函数与弱范数，训练端点的定量判断。
\Needspace{6\baselineskip}
\begin{exercise}\label{b2:ex:18-new-1}
在带 Lebesgue 测度的 \(\mathbb R\) 上，令
\[
 h=H\mathbf1_{(0,1)},\quad
 h(x)=\frac1\pi\log\left|\frac{x}{x-1}\right|\quad(x\ne0,1),
\]
并在两个端点任意补有限值。对 \(t>0\)，求出 \(\{|h|>t\}\) 的区间组成，证明其分布函数与弱 \(L^1\) 准范数分别为
\[
 d_h(t)=\frac2{\sinh(\pi t)},
 \quad \|h\|_{L^{1,\infty}}
       \coloneqq\sup_{t>0}t\cdot d_h(t)=\frac2\pi.
\]
分别解释 \(t\downarrow0\) 和 \(t\uparrow\infty\) 的分布行为，并据此判断 \(h\) 的有限指数可积性。
\end{exercise}
\begin{solution}
固定 \(t>0\)，记 \(q=e^{\pi t}>1\)。条件 \(|h(x)|>t\) 等价于
\[
 \left|\frac{x}{x-1}\right|<q^{-1}
 \quad\text{或}\quad
 \left|\frac{x}{x-1}\right|>q.
\]
在 \(x<0\) 上，比值的绝对值为 \(-x/(1-x)<1\)，第一项不等式给出
\(-1/(q-1)<x<0\)。在 \(0<x<1\) 上，比值为 \(x/(1-x)\)，两项分别给出
\(0<x<1/(q+1)\) 和 \(q/(q+1)<x<1\)。在 \(x>1\) 上，比值为 \(x/(x-1)>1\)，第二项给出
\(1<x<q/(q-1)\)。因此，忽略有限多个点后，超水平集就是
\[
 \left(-\frac1{q-1},\frac1{q+1}\right)
 \ \cup\
 \left(\frac q{q+1},\frac q{q-1}\right).
\]
两段长度之和为
\[
 d_h(t)=2\left(\frac1{q-1}+\frac1{q+1}\right)
       =\frac{4q}{q^2-1}=\frac2{\sinh(\pi t)}.
\]
由 \(\sinh s\geq s\)（\(s>0\)），有 \(t\cdot d_h(t)\leq2/\pi\)；又由
\(\sinh s/s\to1\)（\(s\downarrow0\)），这个上界在极限中达到。所以弱范数等于 \(2/\pi\)，其上确界无需在某个正高度取到。

更具体地，低高度时 \(d_h(t)\sim2/(\pi t)\)，记录的是远处 \(1/x\) 尾部占据的长区间；高高度时 \(d_h(t)\sim4e^{-\pi t}\)，记录的是两个对数奇点附近迅速缩小的区域。用分布函数恒等式
\[
 \|h\|_p^p=p\int_0^\infty t^{p-1}\cdot d_h(t)\,\mathrm dt
 \quad(0<p<\infty),
\]
无穷端因指数衰减总可积，零端则等价于 \(\int_0^1t^{p-2}\,\mathrm dt\)。故恰在 \(p>1\) 时 \(h\in L^p\)，在 \(0<p\leq1\) 时不属于 \(L^p\)。此外高高度的分布函数始终为正，故 \(h\notin L^\infty\)。
\end{solution}

% 自拟：将 Hilbert 变换的代数恒等式与正负频率结合，判断可逆性及精确平方范数。
\Needspace{6\baselineskip}
\begin{exercise}\label{b2:ex:18-new-2}
设 \(1<p<\infty\)，\(a,b\in\mathbb C\)，\(H\) 取本章兼容的 \(L^p(\mathbb R;\mathbb C)\) 延拓。证明 \(A=aI+bH\) 为可逆有界算子，当且仅当 \(a^2+b^2\ne0\)，并在可逆时求其逆。再在 \(L^2\) 上证明
\[
 \|A\|_{2\to2}=\max\{|a-\mathrm ib|,|a+\mathrm ib|\},
 \quad
 \|A^{-1}\|_{2\to2}
 =\frac1{\min\{|a-\mathrm ib|,|a+\mathrm ib|\}}.
\]
必要性与范数下界须给出实际输入，不能只在一个频率点上检验符号。
\end{exercise}
\begin{solution}
先将正文的平方恒等式传到所用空间。在 \(L^2\) 中已有 \(H^2=-I\)。对 \(\varphi\in\mathcal S\)，其 \(L^2\) 与 \(L^p\) 实现相容，且 \(H\varphi\in L^2\cap L^p\)，所以该恒等式也在 \(L^p\) 中成立。对一般 \(f\in L^p\)，取 \(\varphi_n\in\mathcal S\) 在 \(L^p\) 中逼近；两次使用 \(H\) 的有界性，就得到 \(H^2f=-f\)。因此
\[
 (aI+bH)(aI-bH)=(aI-bH)(aI+bH)=(a^2+b^2)I.
\]
若 \(a^2+b^2\ne0\)，则
\[
 A^{-1}=\frac{aI-bH}{a^2+b^2}
\]
是所需的有界双侧逆。

由于 \(a^2+b^2=(a-\mathrm ib)(a+\mathrm ib)\)，退化时至少有一个因子为零。取非零 \(\psi_+\in C_c^\infty\bigl((1,2)\bigr)\)，令 \(f_+=\mathcal F^{-1}\psi_+\in\mathcal S\)。其频率全为正，故
\(Hf_+=-\mathrm if_+\)，从而 \(Af_+=(a-\mathrm ib)f_+\)。同样取非零频率测试函数支撑于 \((-2,-1)\)，得到 \(f_-\in\mathcal S\) 且
\(Af_-=(a+\mathrm ib)f_-\)。每个 Schwartz 函数都属于当前的 \(L^p\)，所以任一因子为零都使 \(A\) 有非零核，不能可逆。

在 \(L^2\) 中，Plancherel 给出
\[
 \|Af\|_2^2
 =|a-\mathrm ib|^2\int_{\xi>0}|\widehat f(\xi)|^2\,\mathrm d\xi
  +|a+\mathrm ib|^2\int_{\xi<0}|\widehat f(\xi)|^2\,\mathrm d\xi.
\]
这给出最大值所对应的上界，而选取 \(f_+\) 或 \(f_-\) 就得到相应下界。可逆时，逆算子在两个半轴上的符号分别为
\((a-\mathrm ib)^{-1}\) 和 \((a+\mathrm ib)^{-1}\)，用同样的输入得到其精确范数。正、负半轴都是正测度集，所以这些范数由真正的 \(L^2\) 函数实现。
\end{solution}

% 自拟：区分算子范数最优与非零函数能否取等，训练频率集中构造。
\Needspace{6\baselineskip}
\begin{exercise}\label{b2:ex:18-new-3}
设整数 \(d\geq2\)、\(1\leq j\leq d\)，\(R_j\) 为乘子
\(-\mathrm i\xi_j/|\xi|\) 定义的 \(L^2(\mathbb R^d;\mathbb C)\) Riesz 变换。证明
\(\|R_j\|_{2\to2}=1\)，但每个非零 \(f\in L^2\) 都满足
\(\|R_jf\|_2<\|f\|_2\)。构造范数为一的 Schwartz 函数族，使其输出范数趋于一。
\end{exercise}
\begin{solution}
由 Plancherel，
\[
 \|f\|_2^2-\|R_jf\|_2^2
 =\int_{\mathbb R^d}
   \frac{\sum_{k\ne j}\xi_k^2}{|\xi|^2}
       \cdot|\widehat f(\xi)|^2\,\mathrm d\xi\geq0.
\]
权重只在第 \(j\) 条坐标轴上为零，原点处的补值不影响积分。由于 \(d\geq2\)，这条轴的 \(d\) 维 Lebesgue 测度为零。若 \(f\ne0\)，则 \(|\widehat f|>0\) 的集合有正测度，故上式严格为正。特别地，范数等式不能由某个非零输入取到。

为逼近上界，记 \(e_j\) 为第 \(j\) 个单位坐标向量，取非零
\(\psi\in C_c^\infty\bigl(B(0,1)\bigr)\)。对 \(0<\varepsilon<1/2\)，定义
\[
 \widehat f_\varepsilon(\xi)
 =\frac{\varepsilon^{-d/2}}{\|\psi\|_2}
       \cdot\psi\left(\frac{\xi-e_j}{\varepsilon}\right).
\]
该函数光滑且紧支撑，所以 \(f_\varepsilon\in\mathcal S\)；换元给出
\(\|f_\varepsilon\|_2=1\)。在其频率支集上，\(\xi_j\geq1-\varepsilon\)、\(|\xi|\leq1+\varepsilon\)，所以
\[
 \left(\frac{1-\varepsilon}{1+\varepsilon}\right)^2
 \leq\|R_jf_\varepsilon\|_2^2\leq1.
\]
令 \(\varepsilon\downarrow0\)，输出范数便趋于一。这说明最优算子范数可以由集中族逼近，即使没有任何非零取等函数。维数条件也有实质作用：在一维，\(R_1=H\) 在全部 \(L^2\) 上等距。
\end{solution}

% 自拟：将坏块零平均推广到有限阶矩消失，并用有限差分构造精确阶数。
\Needspace{6\baselineskip}
\begin{exercise}\label{b2:ex:18-new-4}
设 \(m\geq0\) 为整数、\(R>0\)，\(b\in L^1(\mathbb R;\mathbb C)\) 在 \([-R,R]\) 外几乎处处为零。定义
\[
 \mu_\ell\coloneqq\int_{\mathbb R}y^\ell\cdot b(y)\,\mathrm dy
 \quad(0\leq\ell\leq m),
\]
并假设 \(\mu_\ell=0\) 对 \(0\leq\ell<m\) 成立。当 \(m=0\) 时这一组消失条件为空。此题只在支集外定义普通绝对积分
\[
 Hb(x)\coloneqq\frac1\pi\int_{\mathbb R}\frac{b(y)}{x-y}\,\mathrm dy
 \quad(|x|>R),
\]
不预设一般 \(L^1\) 输入在支集内的主值。证明当 \(|x|>2R\) 时，
\[
 \left|Hb(x)-\frac{\mu_m}{\pi x^{m+1}}\right|
 \leq\frac{2R^{m+1}}{\pi|x|^{m+2}}\|b\|_1.
\]
再对
\[
 b_m(y)=\sum_{k=0}^m(-1)^k\binom mk\cdot\mathbf1_{(k,k+1)}(y)
\]
证明前 \(m\) 个矩为零，而 \(\mu_m=(-1)^m m!\)，从而该尾部阶数可以实际出现。
\end{exercise}
\begin{solution}
紧支集保证每个所用的矩都绝对收敛。在 \(|x|>2R\)、\(|y|\leq R\) 时，有限几何展开给出恒等式
\[
 \frac1{x-y}=\sum_{\ell=0}^m\frac{y^\ell}{x^{\ell+1}}
                  +\frac{y^{m+1}}{x^{m+1}(x-y)}.
\]
乘 \(b(y)\) 积分，前 \(m\) 项因矩条件消失。余项中
\(|x-y|\geq|x|/2\)，故
\[
 \begin{aligned}
 \left|Hb(x)-\frac{\mu_m}{\pi x^{m+1}}\right|
 &\leq\frac2{\pi|x|^{m+2}}
             \int|y|^{m+1}\cdot|b(y)|\,\mathrm dy\\
 &\leq\frac{2R^{m+1}}{\pi|x|^{m+2}}\|b\|_1.
 \end{aligned}
\]
估计对正、负两侧的 \(x\) 都成立，首项的符号保留在 \(x^{m+1}\) 中。

对给定的 \(b_m\)，它支撑于 \([0,m+1]\)，因而属于 \(L^1\)。将每个区间换元为 \(y=z+k\)，得到
\[
 \int y^\ell\cdot b_m(y)\,\mathrm dy
 =\int_0^1\sum_{k=0}^m(-1)^k\binom mk(z+k)^\ell\,\mathrm dz.
\]
记前向差分为 \(\Delta P(z)=P(z+1)-P(z)\)。反复展开可得
\[
 \Delta^mP(z)=\sum_{k=0}^m(-1)^{m-k}\binom mk\cdot P(z+k).
\]
差分每次将多项式次数降低一，且把最高次系数乘以原来的次数。因此对 \(P(z)=z^\ell\)，\(\ell<m\) 时 \(\Delta^mP=0\)，\(\ell=m\) 时 \(\Delta^mP=m!\)。上面的矩积分于是分别为零和 \((-1)^m m!\)；\(m=0\) 时同样得到 \(\mu_0=1\)。由于第 \(m\) 个矩非零，所得尾部的首项确为 \(x^{-(m+1)}\) 阶，而非仅有一个可能不尖锐的上界。
\end{solution}

% 自拟：从离支集局部性识别对角线自由度，补足零核乘法例子的结构解释。
\Needspace{6\baselineskip}
\begin{exercise}\label{b2:ex:18-new-5}
设整数 \(d\geq1\)，\(T_1,T_2\) 是 \(L^2(\mathbb R^d;\mathbb C)\) 上的有界复线性算子。设可测核 \(K\) 在 \(x\ne y\) 上满足
\(|K(x,y)|\leq A|x-y|^{-d}\)，其中 \(A<\infty\)；并且对每个紧集 \(F\) 与每个在 \(F\) 外几乎处处为零的 \(f\in L^2\)，有
\[
 T_1f(x)=T_2f(x)=\int_FK(x,y)\cdot f(y)\,\mathrm dy
 \quad\text{对几乎处处的 }x\notin F.
\]
证明差算子 \(S=T_1-T_2\) 必为乘法算子：存在唯一的 \(a\in L^\infty(\mathbb R^d)\) 等价类，使
\[
 Sf=a\cdot f\quad(f\in L^2),
 \quad\|S\|_{2\to2}=\|a\|_\infty.
\]
证明中须从紧支集输入出发，解释为什么 \(S\) 对任意可测集合都保持支撑，而不能只对开集作形式判断。
\end{exercise}
\begin{solution}
若 \(f\) 支撑于紧集 \(F\)，共同核表示相减便知 \(Sf=0\) 在 \(F\) 外几乎处处成立。先把这个局部性推广到任意 Lebesgue 可测集 \(E\)。设 \(f\in L^2\) 在 \(E\) 外几乎处处为零，选定有限值可测代表，并令
\[
 E_n=E\cap B(0,n)\cap\{|f|\leq n\}.
\]
这些集合递增且覆盖 \(f\) 的非零部分，至多相差零测集。由有限测度集合的内正则性，选紧集 \(K_n\subset E_n\)，使
\(|E_n\setminus K_n|<n^{-4}\)。令 \(f_n=f\cdot\mathbf1_{K_n}\)，则
\[
 \|f-f_n\|_2^2
 \leq\int_{E\setminus E_n}|f|^2+n^2|E_n\setminus K_n|
 \leq\int_{E\setminus E_n}|f|^2+n^{-2}\longrightarrow0.
\]
每个 \(Sf_n\) 都在 \(K_n\subset E\) 外为零，且 \(Sf_n\to Sf\) 于 \(L^2\)。乘以 \(\mathbf1_{E^c}\) 是有界算子，故极限 \(Sf\) 也在 \(E\) 外为零。这里不要求 \(K_n\) 递增。

对任意 \(f\in L^2\)，将它分成 \(\mathbf1_E\cdot f\) 与 \(\mathbf1_{E^c}\cdot f\)，两部分的输出分别留在 \(E\) 与 \(E^c\)。因此
\[
 S(\mathbf1_E\cdot f)=\mathbf1_E\cdot Sf.
\]
现在令 \(D_n=(-n,n)^d\)、\(a_n=S\mathbf1_{D_n}\)。当 \(m<n\) 时，
\[
 \mathbf1_{D_m}\cdot a_n
 =S(\mathbf1_{D_m}\cdot\mathbf1_{D_n})
 =S\mathbf1_{D_m}=a_m.
\]
取去掉可数个零测例外集的共同满测集，这些函数就在各 \(D_n\) 上相容，可以拼成一个可测函数 \(a\)，满足 \(a=a_n\) 于 \(D_n\)。对每个可测 \(E\subset D_n\)，又有
\[
 S\mathbf1_E=\mathbf1_E\cdot a,
 \quad\|a\cdot\mathbf1_E\|_2
      \leq\|S\|_{2\to2}\cdot|E|^{1/2}.
\]
若 \(\{|a|>\|S\|_{2\to2}+\varepsilon\}\cap D_n\) 有正测度，以上式中的 \(E\) 取为此集合，平方后就与 \(L^2\) 积分下界矛盾。因此对每个 \(n\) 和 \(\varepsilon>0\)，该集合测度为零，得到 \(\|a\|_\infty\leq\|S\|_{2\to2}\)。

由线性，\(Sf=a\cdot f\) 对支撑在某个 \(D_n\) 内的简单函数成立。这样的简单函数在 \(L^2\) 中稠密：先截去空间远尾和过高幅度，再将有界的复值范围作有限网格近似。两边均为有界算子，故等式推广到全部 \(L^2\)。反向不等式
\(\|S\|_{2\to2}=\|M_a\|_{2\to2}\leq\|a\|_\infty\)
来自直接积分，于是范数相等。若另一个 \(\widetilde a\) 给出相同乘法算子，作用于各 \(\mathbf1_{D_n}\) 即得 \(a=\widetilde a\) 几乎处处，证明唯一性。可见离支集核所遗漏的，恰好是这样的有界乘法项。
\end{solution}

% 自拟：组合方向乘子得到向量投影，连接奇异积分与无散条件。
\Needspace{6\baselineskip}
\begin{exercise}\label{b2:ex:18-new-6}
设整数 \(d\geq2\)、\(1<p<\infty\)，\(R_j\) 取本章兼容的 \(L^p(\mathbb R^d;\mathbb C)\) Riesz 变换。对向量函数 \(F=(F_1,\ldots,F_d)\in L^p(\mathbb R^d;\mathbb C^d)\)，定义
\[
 (\mathbb PF)_i=F_i+\sum_{j=1}^dR_iR_jF_j
 \quad(1\leq i\leq d),
\]
其中向量范数采用逐点欧氏范数。证明 \(\mathbb P\) 在该向量值 \(L^p\) 空间上有界，\(\mathbb P^2=\mathbb P\)，且
\(\operatorname{div}(\mathbb PF)=0\) 于分布意义；对每个 \(\varphi\in\mathcal S(\mathbb R^d)\)，还有 \(\mathbb P\nabla\varphi=0\)。再证明在 \(p=2\) 时，\(\mathbb P\) 是到闭子空间
\[
 \mathcal V=\{F\in L^2(\mathbb R^d;\mathbb C^d):
                         \operatorname{div}F=0\text{ 于分布意义}\}
\]
的正交投影，且 \(\|\mathbb P\|_{2\to2}=1\)。
\end{exercise}
\begin{solution}
先看有界性。每个分量的 \(L^p\) 范数不超过 \(\|F\|_{L^p(\mathbb C^d)}\)，且 \(R_i,R_j\) 均为有界 \(L^p\) 算子。有限次使用三角不等式，例如
\[
 \begin{aligned}
 \|\mathbb PF\|_{L^p(\mathbb C^d)}
 &\leq\sum_{i=1}^d\|(\mathbb PF)_i\|_p\\
 &\leq\sum_i\|F_i\|_p
   +\sum_{i,j}\|R_i\|_{p\to p}\cdot\|R_j\|_{p\to p}\cdot\|F_j\|_p\\
 &\leq C_{d,p}\|F\|_{L^p(\mathbb C^d)},
 \end{aligned}
\]
就给出足够的界，不要求这个常数最优。

在 \(L^2\) 中，\(R_iR_j\) 的符号为 \(-\xi_i\xi_j/|\xi|^2\)，所以 \(\mathbb P\) 的矩阵符号是
\[
 P(\xi)=I-\frac{\xi\xi^{\mathsf T}}{|\xi|^2}
 \quad(\xi\ne0).
\]
它是到 \(\xi^\perp\subset\mathbb C^d\) 的欧氏正交投影，因此
\(P(\xi)^2=P(\xi)=P(\xi)^*\)、\(\xi^{\mathsf T}P(\xi)=0\)，且矩阵范数不超过一。Plancherel 把这些等式传给 \(\mathbb P\)，得到幂等性、自伴性、平方范数上界及无散性。

还须识别投影的像。对 \(F\in L^2\)，分布恒等式 \(\operatorname{div}F=0\) 等价于
\(\xi\cdot\widehat F(\xi)=0\) 几乎处处：Fourier 变换将散度化为 \(2\pi\mathrm i\xi\cdot\widehat F\)，后者局部可积，一个局部可积函数给出零分布当且仅当它几乎处处为零。这样的 \(F\) 满足 \(P(\xi)\widehat F=\widehat F\)，即 \(\mathbb PF=F\)。而任意输出已经无散，所以像恰为 \(\mathcal V\)。若一列无散向量在 \(L^2\) 中收敛，与每个紧支集光滑测试函数的梯度配对都收敛，极限仍无散，故 \(\mathcal V\) 闭。结合自伴与幂等，得到所述正交投影。

范数确实为一，而不只是至多一。取 \(\psi\in C_c^\infty(\mathbb R^d)\)，使
\((\partial_1\psi,\partial_2\psi)\) 不恒为零，并令
\[
 G=(\partial_2\psi,-\partial_1\psi,0,\ldots,0).
\]
混合偏导相等使 \(\operatorname{div}G=0\)，且 \(G\ne0\)，故 \(\mathbb PG=G\) 给出下界一。另一方面，\(\widehat{\nabla\varphi}=2\pi\mathrm i\xi\cdot\widehat\varphi\) 沿 \(\xi\) 方向，矩阵 \(P(\xi)\) 将其消去，所以 \(\mathbb P\nabla\varphi=0\) 在 \(L^2\) 中成立。

最后处理一般 \(p\)。对 Schwartz 向量，上述等式先在 \(L^2\) 成立；它和其中的 \(\mathbb P\) 输出均属于 \(L^2\cap L^p\)，各个 Riesz 变换的相容性使幂等式也在 \(L^p\) 中成立。取 Schwartz 向量在 \(L^p\) 中逼近任意 \(F\)，利用有界性即可传递 \(\mathbb P^2=\mathbb P\)。对每个紧支集光滑测试函数，其梯度属于 \(L^{p'}\)，所以输出的无散分布等式也能在该极限中传递。梯度消去公式同样由相容性成立。这里的正交性与范数一专指 \(p=2\)，未对其他指数作相同断言。
\end{solution}

% 自拟：把热核积分、预解方程与二阶 Riesz 乘子连接，并区分有质量与零质量估计。
\Needspace{6\baselineskip}
\begin{exercise}\label{b2:ex:18-new-7}
设整数 \(d\geq1\)、\(1<p<\infty\)、\(\lambda>0\)，\(f\in L^p(\mathbb R^d;\mathbb C)\)。令
\[
 H_t(x)=(4\pi t)^{-d/2}e^{-|x|^2/(4t)},
 \quad u_\lambda=\int_0^\infty e^{-\lambda t}\cdot(H_t*f)\,\mathrm dt,
\]
时间积分按 \(L^p\) 范数中的积分理解。证明它收敛，且
\[
 (\lambda-\Delta)u_\lambda=f\text{ 于分布意义},
 \quad\|u_\lambda\|_p\leq\lambda^{-1}\|f\|_p.
\]
先对 \(f\in\mathcal S\) 求出 \(\widehat u_\lambda\) 并证明 \(u_\lambda\in\mathcal S\)，再由
\[
 \partial_i\partial_ju_\lambda
 =R_iR_j(f-\lambda u_\lambda)
\]
证明全部二阶分布导数属于 \(L^p\)，且
\[
 \sum_{i,j=1}^d\|\partial_i\partial_ju_\lambda\|_p
 \leq C_{d,p}\|f\|_p,
\]
其中常数与 \(\lambda>0\) 无关。最后用缩放说明不存在对所有 \(u\in C_c^\infty(\mathbb R^d)\) 都成立的统一估计
\(\|u\|_p\leq C\|\Delta u\|_p\)。
\end{exercise}
\begin{solution}
先核对时间积分的连续性。固定 \(t_0>0\)，当 \(t_0/2\leq t\leq2t_0\) 时，热核的显式公式给出共同可积控制 \(C_{d,t_0}e^{-|x|^2/(8t_0)}\)。支配收敛说明 \(\|H_t-H_{t_0}\|_1\to0\)，Young 不等式因而给出
\[
 \|H_t*f-H_{t_0}*f\|_p
 \leq\|H_t-H_{t_0}\|_1\cdot\|f\|_p\longrightarrow0.
\]
零时的 \(L^p\) 连续性由\cref{b2:thm:12-heat-convolution}给出。热核非负且积分为一，Young 不等式还给出
\[
 \|H_t*f\|_p\leq\|f\|_p,
 \quad\int_0^\infty e^{-\lambda t}\|H_t*f\|_p\,\mathrm dt
                  \leq\lambda^{-1}\|f\|_p.
\]
因而可先在紧时间区间内积分，再放开两端，得到收敛的 \(L^p\) 值积分及所求零阶界。特别地，\(f\mapsto u_\lambda\) 是范数不超过 \(\lambda^{-1}\) 的线性算子。

若 \(f\in\mathcal S\)，还可在 \(L^1\) 中作同一积分，因为
\(\int_0^\infty e^{-\lambda t}\|H_t*f\|_1\,\mathrm dt<\infty\)。Fubini 因而允许先取 Fourier 变换，得到
\[
 \widehat u_\lambda(\xi)
 =\int_0^\infty e^{-t(\lambda+4\pi^2|\xi|^2)}
                 \cdot\widehat f(\xi)\,\mathrm dt
 =\frac{\widehat f(\xi)}{\lambda+4\pi^2|\xi|^2}.
\]
分母严格为正，其倒数是光滑函数，各阶导数至多作多项式增长，所以乘上 Schwartz 函数仍为 Schwartz 函数。Fourier 反变换说明 \(u_\lambda\in\mathcal S\)，并立即给出
\((\lambda-\Delta)u_\lambda=f\)。

记 \(g=f-\lambda u_\lambda\)。由刚求出的公式，
\[
 \widehat g(\xi)=\frac{4\pi^2|\xi|^2}{\lambda+4\pi^2|\xi|^2}
                       \cdot\widehat f(\xi).
\]
两个 Riesz 符号相乘为 \(-\xi_i\xi_j/|\xi|^2\)，故
\[
 \widehat{R_iR_jg}(\xi)
 =-\frac{4\pi^2\xi_i\xi_j}{\lambda+4\pi^2|\xi|^2}
                       \cdot\widehat f(\xi)
 =\widehat{\partial_i\partial_ju_\lambda}(\xi).
\]
零频处任意补值不影响等式。使用 \(R_iR_j\) 的 \(L^p\) 有界性与零阶界，有
\[
 \begin{aligned}
 \sum_{i,j}\|\partial_i\partial_ju_\lambda\|_p
 &\leq C_{d,p}\|f-\lambda u_\lambda\|_p\\
 &\leq C_{d,p}(\|f\|_p+\lambda\|u_\lambda\|_p)
 \leq2C_{d,p}\|f\|_p.
 \end{aligned}
\]
吸收因子二即可得到与 \(\lambda\) 无关的常数。

对一般 \(f\in L^p\)，取 \(f_n\in\mathcal S\) 在 \(L^p\) 中逼近。零阶界使对应的 \(u_{\lambda,n}\to u_\lambda\) 于 \(L^p\)，而
\(R_iR_j(f_n-\lambda u_{\lambda,n})\)
也在 \(L^p\) 中收敛到 \(R_iR_j(f-\lambda u_\lambda)\)。对任意 \(\varphi\in C_c^\infty\)，Hölder 不等式保证
\[
 \int u_{\lambda,n}\cdot\partial_i\partial_j\varphi
 \longrightarrow\int u_\lambda\cdot\partial_i\partial_j\varphi.
\]
所以二阶分布导数的极限正是 \(\partial_i\partial_ju_\lambda\)，并由上述 \(L^p\) 极限表示。对方程作同样的测试配对，得到其分布版本；范数估计也随 \(L^p\) 收敛传递。

零质量参数不能提供同样的零阶控制。固定非零 \(\phi\in C_c^\infty\)，使 \(\Delta\phi\ne0\)，令 \(v_R(x)=\phi(x/R)\)、\(R>0\)。换元给出
\[
 \|v_R\|_p=R^{d/p}\|\phi\|_p,
 \quad\|\Delta v_R\|_p=R^{d/p-2}\|\Delta\phi\|_p.
\]
若存在所述统一常数，则 \(R^2\|\phi\|_p\leq C\|\Delta\phi\|_p\) 对所有 \(R\) 成立，与 \(R\uparrow\infty\) 矛盾。二阶估计中相同的缩放因子可以相消，零阶估计却不能，这解释了 \(\lambda>0\) 在前一估计中的作用。
\end{solution}

% 自拟：用光滑乘法系数的差消去对角线奇性，比较差核的逐点极限与平方范数极限。
\Needspace{6\baselineskip}
\begin{exercise}\label{b2:ex:18-new-8}
设 \(d\geq1\) 为整数，\(T\) 是 \(\mathbb R^d\) 上\cref{b2:thm:18-cz-operator-definition}的标准 Calder\'on--Zygmund 算子，核为 \(K\)。另设硬截断
\[
 T_\epsilon g(x)=\int_{|x-y|>\epsilon}K(x,y)\cdot g(y)\,\mathrm dy
\]
对每个 \(g\in C_c^\infty(\mathbb R^d;\mathbb C)\)，当 \(\epsilon\downarrow0\) 时在 \(L^2\) 中趋于 \(Tg\)。给定 \(a,f\in C_c^\infty(\mathbb R^d;\mathbb C)\)，令 \(C_af=a\cdot Tf-T(a\cdot f)\)。证明
\[
 C_af(x)=\int_{\mathbb R^d}\bigl(a(x)-a(y)\bigr)\cdot K(x,y)\cdot f(y)\,\mathrm dy
\]
几乎处处成立，且右端对每个 \(x\) 都绝对收敛；解释乘法系数的差如何改善对角线附近的奇性。再证明对每个 \(1<p<\infty\)，有
\[
 \|C_af\|_p\leq2\|a\|_\infty\cdot\|T\|_{L^p\to L^p}\cdot\|f\|_p.
\]
这里 \(T\) 的 \(L^p\) 实现取\cref{b2:thm:18-cz-lp}的兼容延拓，不预先假定 \(T_\epsilon f\) 单独几乎处处收敛。
\end{exercise}
\begin{solution}
设核的大小常数为 \(A\)，\(a\) 的全局 Lipschitz 常数为 \(L_a\)。当 \(|x-y|<1\) 时，
\[
 \Bigl|\bigl(a(x)-a(y)\bigr)\cdot K(x,y)\cdot f(y)\Bigr|
 \leq A L_a\|f\|_\infty\cdot|x-y|^{1-d}.
\]
这在单位球上可积，径向积分为 \(\int_0^1\mathrm dr\)；系数之差提供的 \(|x-y|\) 消去了一阶奇性。在 \(|x-y|\geq1\) 时，绝对积分不超过 \(2A\|a\|_\infty\cdot\|f\|_1\)。因此对每个 \(x\)，差核的全积分绝对收敛，并且
\[
 a(x)\cdot T_\epsilon f(x)-T_\epsilon(a\cdot f)(x)
 =\int_{|x-y|>\epsilon}\bigl(a(x)-a(y)\bigr)\cdot K(x,y)\cdot f(y)\,\mathrm dy
\]
随 \(\epsilon\downarrow0\) 趋于该积分。另一方面，\(f,a\cdot f\in C_c^\infty\)，题设的 \(L^2\) 收敛及乘法算子的有界性使左端在 \(L^2\) 中趋于 \(a\cdot Tf-T(a\cdot f)\)。从任意趋零截断序列中抽取一个 \(L^2\) 收敛的逐点子列，就把后者与已存在的全参数差核极限识别，得到所求几乎处处等式。这一步只比较两项之差，没有预设两项各自的逐点极限。

最后利用 \(T\) 的 \(L^p\) 有界性及 \(L^2\)、\(L^p\) 实现的相容性，
\[
 \|a\cdot Tf-T(a\cdot f)\|_p
 \leq\|a\|_\infty\cdot\|Tf\|_p+\|T(a\cdot f)\|_p
 \leq2\|a\|_\infty\cdot\|T\|_{L^p\to L^p}\cdot\|f\|_p.
\]
绝对积分表示使用了 \(a\) 的 Lipschitz 正则性，而这个范数估计只用其有界性，两项结论的假设作用不同。
\end{solution}
